% Generated draft: requires mathematical and source review.
\documentclass[11pt]{article}
\usepackage[margin=25mm]{geometry}
\usepackage{fontspec}
\setmainfont{Latin Modern Roman}
\usepackage{amsmath,amssymb,mathtools,unicode-math}
\setmathfont{Latin Modern Math}
\usepackage{enumitem,needspace,xurl,hyperref}
\setlength{\parindent}{0pt}
\setlength{\parskip}{5pt}
\setlength{\emergencystretch}{3em}
\newcommand{\N}{\mathbb{N}}
\newcommand{\Z}{\mathbb{Z}}
\newcommand{\Q}{\mathbb{Q}}
\newcommand{\R}{\mathbb{R}}
\newcommand{\C}{\mathbb{C}}
\newcommand{\F}{\mathbb{F}}
\newcommand{\A}{\mathbb{A}}
\newcommand{\PP}{\mathbb{P}}
\newcommand{\E}{\mathbb{E}}
\newcommand{\eps}{\varepsilon}
\newcommand{\dd}{\,\mathrm{d}}
\begin{document}
\section*{UnsolvedMath ID: 30002326}
\textbf{Source problem code:} catalogue entry
\subsection{Definitions and mathematical statement}
For a finite connected graph \(G\), let \(V(G)\) be its vertex set and let \(d_G(u,v)\) denote the length of a shortest path between \(u\) and \(v\) in \(G\). For \(S\subseteq V(G)\), define its multiset of pairwise distances by
\[
D_G(S)=\bigl\{\!\bigl\{d_G(u,v):\{u,v\}\subseteq S,\ u\ne v\bigr\}\!\bigr\},
\]
where each unordered pair of distinct vertices contributes once, and equal distances retain their multiplicities. Two subsets \(A,B\subseteq V(G)\) are \emph{homometric} if \(D_G(A)=D_G(B)\); they are disjoint if \(A\cap B=\varnothing\).

The supplied statement calls \(h(n)\) the largest integer with a universal existence property, but does not explicitly connect that integer to the cardinalities of the subsets. The following precise formulation is therefore conditional on the intended requirement being \(|A|=|B|=h(n)\). It also explicitly adopts finite simple undirected unweighted graphs and distances measured in the ambient graph \(G\); these conventions are not expressly specified in the supplied statement.

Under these conventions, for each integer \(n\ge 2\), define
\[
\begin{aligned}
h(n)=\max\bigl\{k\in\{1,\ldots,\lfloor n/2\rfloor\}:{}&\text{ for every connected graph }G\text{ with }|V(G)|=n,\\
&\text{ there exist }A,B\subseteq V(G)\text{ with }A\cap B=\varnothing,\\
&|A|=|B|=k\text{ and }D_G(A)=D_G(B)\bigr\}.
\end{aligned}
\]
The question is whether
\[
h(n)=o(n)\qquad(n\to\infty),
\]
that is, whether
\[
\forall\varepsilon>0\ \exists N\ge 2\ \forall n\in\Z,\qquad n\ge N\Longrightarrow h(n)<\varepsilon n.
\]
Restricting this formulation to \(n\ge2\) avoids imposing an unspecified convention concerning empty subsets at \(n=1\) and suffices for the asymptotic question. The missing cardinality requirement and the graph and distance conventions require confirmation against the cited source before this conditional formulation is accepted.
\subsection{Short English statement}
Does the largest size of two disjoint homometric vertex subsets guaranteed in every connected \(n\)-vertex graph grow sublinearly in \(n\)? Here homometric means having identical multisets of pairwise distances. This draft makes the intended subset size and graph conventions explicit, subject to source confirmation.
\subsection{Sources}
\begin{itemize}
\item Local UnsolvedMath registry, supplied complete source record: permanent ID \(30002326\), title \emph{Sublinear Bounds for Disjoint Homometric Vertex Sets}, exact source problem code \texttt{catalogue entry}, in the Oberwolfach Reports Open Problems collection. The record gives proposed year 2013 and does not name a proposer.
\item Supplied report citation: \emph{Combinatorics and Probability (2014)}, Oberwolfach Reports, DOI 10.4171/owr/2013/18, pp.\ 1087--1152. \url{https://doi.org/10.4171/owr/2013/18}. The supplied background describes Oberwolfach Reports as documenting mathematical workshops and their open-problem discussions. The citation's year 2014 differs from the record's proposed year 2013 and the year appearing in the DOI; these supplied details are retained without reconciliation.
\item Supplied literature source: Axenovich--Özkahya, homometric sets. \url{https://arxiv.org/abs/1203.1158}.
\item Supplied literature source: Oberwolfach combinatorics report. \url{https://doi.org/10.4171/owr/2013/18}.
\item Literature assessment dated exactly \texttt{2026-08-21}: the supplied curation labels the problem open and states that known bounds for the homometric-set parameter remain far from proving \(h(n)=o(n)\). This is supplied dated curation derived from the OWR record, not an independently verified assessment; time-sensitive status and cited claims require independent verification.
\item The supplied record labels the statement \texttt{corrected\_verified} and reports that its mathematical content was checked against the official Oberwolfach report, with notation, context, and question form normalized without changing the source-level claim. This is a supplied verification assertion; no independent source verification was performed for this draft.
\end{itemize}
\end{document}
